
\documentclass{article}
\usepackage{colortbl}
\usepackage{makecell}
\usepackage{multirow}
\usepackage{supertabular}

\begin{document}

\newcounter{utterance}

\centering \large Interaction Transcript for game `clean\_up', experiment `2\_hard\_3obj\_nl', episode 2 with gemma{-}4{-}26b{-}a4b{-}it.
\vspace{24pt}

{ \footnotesize  \setcounter{utterance}{1}
\setlength{\tabcolsep}{0pt}
\begin{supertabular}{c@{$\;$}|p{.15\linewidth}@{}p{.15\linewidth}p{.15\linewidth}p{.15\linewidth}p{.15\linewidth}p{.15\linewidth}}
    \# & $\;$A & \multicolumn{4}{c}{Game Master} & $\;\:$B\\
    \hline

    \theutterance \stepcounter{utterance}  
    & & \multicolumn{4}{p{0.6\linewidth}}{
        \cellcolor[rgb]{0.9,0.9,0.9}{
            \makecell[{{p{\linewidth}}}]{
                \texttt{\tiny{[P1$\langle$GM]}}
                \texttt{Ik ben uw spelleider, en u speelt een samenwerkingsspel met het volgende raster als speelbord:} \\
\\ 
\texttt{```} \\
\texttt{ 1234567} \\
\texttt{╔═╤══╤══╗} \\
\texttt{║◌│◌◌│◌◌║ 1} \\
\texttt{║◌└┬─┴──╢ 2} \\
\texttt{║◌◌│L◌◌◌║ 3} \\
\texttt{╟─┬┘◌◌┌─╢ 4} \\
\texttt{║◌│◌◌◌│C║ 5} \\
\texttt{╟─┼─┬─┤P║ 6} \\
\texttt{║◌│◌│◌│◌║ 7} \\
\texttt{╚═╧═╧═╧═╝} \\
\\ 
\texttt{```} \\
\\ 
\texttt{* De bovenrand toont x{-}coördinaten die naar rechts oplopen, en de rechterrand y{-}coördinaten die naar beneden oplopen.} \\
\texttt{* De volgende objecten zijn willekeurig op uw raster geplaatst: 'C', 'L', 'P'.} \\
\\ 
\\ 
\texttt{De andere speler ziet een variatie van het speelbord, waarbij de objecten op verschillende willekeurige locaties zijn geplaatst. U kunt het bord van de andere speler niet zien en zij kunnen het uwe niet zien.} \\
\\ 
\\ 
\texttt{**Doel:**} \\
\\ 
\texttt{Beide spelers moeten de objecten op hun respectievelijke achtergrond zo verplaatsen dat identieke objecten op dezelfde coördinaten terechtkomen. U moet communiceren met de andere speler om een gemeenschappelijke doelconfiguratie overeen te komen.} \\
\\ 
\\ 
\texttt{**Regels:**} \\
\\ 
\texttt{* In elke beurt kunt u precies één van de volgende twee commando's sturen:} \\
\texttt{ 1. `ZEG: <BERICHT>`: om een bericht (alles tot de volgende regeleinde) naar de andere speler te sturen. Ik zal het doorsturen naar uw partner.} \\
\texttt{2. `BEWEEG: <OBJECT>, (<X>, <Y>)`: om een object naar een nieuwe positie te verplaatsen, waarbij `<X>` de kolom is en `<Y>` de rij. Ik zal u laten weten of uw zet succesvol was of niet.} \\
\texttt{* Als u zich niet aan het formaat houdt, of meerdere commando's tegelijk verstuurt, moet ik u straffen.} \\
\texttt{* Als beide spelers meer dan 8 strafpunten verzamelen, verliezen jullie allebei het spel.} \\
\texttt{* Het is van vitaal belang dat u met de andere speler communiceert over uw doelstatus! De *enige* manier waarop u uw strategie aan de andere speler kunt overbrengen is door gebruik te maken van het `ZEG: <BERICHT>` commando!} \\
\\ 
\\ 
\texttt{**Objecten verplaatsen**} \\
\\ 
\texttt{* U kunt alleen objecten verplaatsen naar cellen binnen de grenzen van het raster. De doelcel moet leeg zijn, d.w.z. deze mag alleen het symbool '◌' bevatten.} \\
\texttt{* Als u probeert een object te verplaatsen naar een plek die niet leeg is, of als u probeert het buiten het raster te verplaatsen, moet ik u straffen. U krijgt een nieuwe poging.} \\
\texttt{* Voordat u een zet doet, moet u controleren of de doelplek leeg is, en geen letter, kader of lijn bevat!} \\
\\ 
\\ 
\texttt{**Einde van het spel**} \\
\\ 
\texttt{Als u denkt dat u het doel bereikt hebt om alle objecten uit te lijnen, kunt u de andere speler vragen het spel te beëindigen door `ZEG: klaar?` te sturen. Als de andere speler u vraagt om het spel te beëindigen en u antwoordt met `ZEG: klaar!`, zal het spel eindigen.} \\
\\ 
\texttt{Beide spelers winnen als het spel eindigt binnen 12 rondes, waarbij één ronde gedefinieerd wordt als twee spelers die elk een geldig commando sturen.} \\
\\ 
\texttt{**Score:**} \\
\\ 
\texttt{Hoe dichter de identieke objecten bij elkaar liggen op beide speelborden, hoe meer punten u krijgt. Strafpunten verminderen uw punten. Kunt u het record verslaan?} \\
\\ 
\\ 
\texttt{Stuur een bericht naar de andere speler om het spel te starten!} \\
            }
        }
    }
    & \\ \\

    \theutterance \stepcounter{utterance}  
    & \multicolumn{4}{p{0.6\linewidth}}{
        \cellcolor[rgb]{0.9,0.9,0.9}{
            \makecell[{{p{\linewidth}}}]{
                \texttt{\tiny{[P1$\rangle$GM]}}
                \texttt{ZEG: Hallo! Laten we samenwerken om de objecten C, L en P op dezelfde coördinaten te krijgen. Ik stel voor dat we ze allemaal naar een centrale plek verplaatsen, bijvoorbeeld naar de coördinaten (4, 4). Wat vind jij daarvan? Laat me weten waar jouw objecten nu staan!} \\
            }
        }
    }
    & & \\ \\

    \theutterance \stepcounter{utterance}  
    & & & \multicolumn{2}{p{0.3\linewidth}}{
        \cellcolor[rgb]{0.9,0.9,0.9}{
            \makecell[{{p{\linewidth}}}]{
                \texttt{\tiny{[GM$|$GM]}}
                \texttt{ZEG: Hallo! Laten we samenwerken om de objecten C, L en P op dezelfde coördinaten te krijgen. Ik stel voor dat we ze allemaal naar een centrale plek verplaatsen, bijvoorbeeld naar de coördinaten (4, 4). Wat vind jij daarvan? Laat me weten waar jouw objecten nu staan!} \\
            }
        }
    }
    & & \\ \\

    \theutterance \stepcounter{utterance}  
    & & \multicolumn{4}{p{0.6\linewidth}}{
        \cellcolor[rgb]{0.9,0.9,0.9}{
            \makecell[{{p{\linewidth}}}]{
                \texttt{\tiny{[GM$\rangle$P2]}}
                \texttt{Ik ben uw spelleider, en u speelt een samenwerkingsspel met het volgende raster als speelbord:} \\
\\ 
\texttt{```} \\
\texttt{ 1234567} \\
\texttt{╔═╤══╤══╗} \\
\texttt{║◌│◌◌│◌◌║ 1} \\
\texttt{║L└┬─┴──╢ 2} \\
\texttt{║◌C│◌◌◌◌║ 3} \\
\texttt{╟─┬┘◌◌┌─╢ 4} \\
\texttt{║◌│◌◌◌│◌║ 5} \\
\texttt{╟─┼─┬─┤◌║ 6} \\
\texttt{║◌│P│◌│◌║ 7} \\
\texttt{╚═╧═╧═╧═╝} \\
\\ 
\texttt{```} \\
\\ 
\texttt{* De bovenrand toont x{-}coördinaten die naar rechts oplopen, en de rechterrand y{-}coördinaten die naar beneden oplopen.} \\
\texttt{* De volgende objecten zijn willekeurig op uw raster geplaatst: 'C', 'L', 'P'.} \\
\\ 
\\ 
\texttt{De andere speler ziet een variatie van het speelbord, waarbij de objecten op verschillende willekeurige locaties zijn geplaatst. U kunt het bord van de andere speler niet zien en zij kunnen het uwe niet zien.} \\
\\ 
\\ 
\texttt{**Doel:**} \\
\\ 
\texttt{Beide spelers moeten de objecten op hun respectievelijke achtergrond zo verplaatsen dat identieke objecten op dezelfde coördinaten terechtkomen. U moet communiceren met de andere speler om een gemeenschappelijke doelconfiguratie overeen te komen.} \\
\\ 
\\ 
\texttt{**Regels:**} \\
\\ 
\texttt{* In elke beurt kunt u precies één van de volgende twee commando's sturen:} \\
\texttt{ 1. `ZEG: <BERICHT>`: om een bericht (alles tot de volgende regeleinde) naar de andere speler te sturen. Ik zal het doorsturen naar uw partner.} \\
\texttt{2. `BEWEEG: <OBJECT>, (<X>, <Y>)`: om een object naar een nieuwe positie te verplaatsen, waarbij `<X>` de kolom is en `<Y>` de rij. Ik zal u laten weten of uw zet succesvol was of niet.} \\
\texttt{* Als u zich niet aan het formaat houdt, of meerdere commando's tegelijk verstuurt, moet ik u straffen.} \\
\texttt{* Als beide spelers meer dan 8 strafpunten verzamelen, verliezen jullie allebei het spel.} \\
\texttt{* Het is van vitaal belang dat u met de andere speler communiceert over uw doelstatus! De *enige* manier waarop u uw strategie aan de andere speler kunt overbrengen is door gebruik te maken van het `ZEG: <BERICHT>` commando!} \\
\\ 
\\ 
\texttt{**Objecten verplaatsen**} \\
\\ 
\texttt{* U kunt alleen objecten verplaatsen naar cellen binnen de grenzen van het raster. De doelcel moet leeg zijn, d.w.z. deze mag alleen het symbool '◌' bevatten.} \\
\texttt{* Als u probeert een object te verplaatsen naar een plek die niet leeg is, of als u probeert het buiten het raster te verplaatsen, moet ik u straffen. U krijgt een nieuwe poging.} \\
\texttt{* Voordat u een zet doet, moet u controleren of de doelplek leeg is, en geen letter, kader of lijn bevat!} \\
\\ 
\\ 
\texttt{**Einde van het spel**} \\
\\ 
\texttt{Als u denkt dat u het doel bereikt hebt om alle objecten uit te lijnen, kunt u de andere speler vragen het spel te beëindigen door `ZEG: klaar?` te sturen. Als de andere speler u vraagt om het spel te beëindigen en u antwoordt met `ZEG: klaar!`, zal het spel eindigen.} \\
\\ 
\texttt{Beide spelers winnen als het spel eindigt binnen 12 rondes, waarbij één ronde gedefinieerd wordt als twee spelers die elk een geldig commando sturen.} \\
\\ 
\texttt{**Score:**} \\
\\ 
\texttt{Hoe dichter de identieke objecten bij elkaar liggen op beide speelborden, hoe meer punten u krijgt. Strafpunten verminderen uw punten. Kunt u het record verslaan?} \\
\\ 
\texttt{De andere speler begon het spel door dit bericht te sturen:} \\
\\ 
\texttt{"Hallo! Laten we samenwerken om de objecten C, L en P op dezelfde coördinaten te krijgen. Ik stel voor dat we ze allemaal naar een centrale plek verplaatsen, bijvoorbeeld naar de coördinaten (4, 4). Wat vind jij daarvan? Laat me weten waar jouw objecten nu staan!"} \\
\\ 
\texttt{Wat is uw eerste opdracht?} \\
            }
        }
    }
    & \\ \\

    \theutterance \stepcounter{utterance}  
    & & & \multicolumn{4}{p{0.6\linewidth}}{
        \cellcolor[rgb]{0.9,0.9,0.9}{
            \makecell[{{p{\linewidth}}}]{
                \texttt{\tiny{[GM$\langle$P2]}}
                \texttt{ZEG: Dat is een goed plan! Laten we (4, 4) als doelwit gebruiken. Op mijn bord staan de objecten nu op de volgende posities: L op (1, 2), C op (2, 3) en P op (3, 7).} \\
            }
        }
    }
     \\ \\

    \theutterance \stepcounter{utterance}  
    & & & \multicolumn{2}{p{0.3\linewidth}}{
        \cellcolor[rgb]{0.9,0.9,0.9}{
            \makecell[{{p{\linewidth}}}]{
                \texttt{\tiny{[GM$|$GM]}}
                \texttt{ZEG: Dat is een goed plan! Laten we (4, 4) als doelwit gebruiken. Op mijn bord staan de objecten nu op de volgende posities: L op (1, 2), C op (2, 3) en P op (3, 7).} \\
            }
        }
    }
    & & \\ \\

    \theutterance \stepcounter{utterance}  
    & & \multicolumn{4}{p{0.6\linewidth}}{
        \cellcolor[rgb]{0.9,0.9,0.9}{
            \makecell[{{p{\linewidth}}}]{
                \texttt{\tiny{[P1$\langle$GM]}}
                \texttt{Je bericht is doorgestuurd naar de andere speler.} \\
\\ 
\\ 
\texttt{Je speelt momenteel ronde 1 van maximaal 12.} \\
\texttt{Jullie hebben gezamenlijk 0 van 8 straffen verzameld.} \\
\texttt{De andere speler stuurde dit bericht:} \\
\\ 
\\ 
\texttt{"Dat is een goed plan! Laten we (4, 4) als doelwit gebruiken. Op mijn bord staan de objecten nu op de volgende posities: L op (1, 2), C op (2, 3) en P op (3, 7)."} \\
\\ 
\\ 
\texttt{Wat is je volgende opdracht?} \\
            }
        }
    }
    & \\ \\

    \theutterance \stepcounter{utterance}  
    & \multicolumn{4}{p{0.6\linewidth}}{
        \cellcolor[rgb]{0.9,0.9,0.9}{
            \makecell[{{p{\linewidth}}}]{
                \texttt{\tiny{[P1$\rangle$GM]}}
                \texttt{BEWEEG: L, (4, 4)} \\
            }
        }
    }
    & & \\ \\

    \theutterance \stepcounter{utterance}  
    & & & \multicolumn{2}{p{0.3\linewidth}}{
        \cellcolor[rgb]{0.9,0.9,0.9}{
            \makecell[{{p{\linewidth}}}]{
                \texttt{\tiny{[GM$|$GM]}}
                \texttt{BEWEEG: L, (4, 4)} \\
            }
        }
    }
    & & \\ \\

    \theutterance \stepcounter{utterance}  
    & & & \multicolumn{2}{p{0.3\linewidth}}{
        \cellcolor[rgb]{0.9,0.9,0.9}{
            \makecell[{{p{\linewidth}}}]{
                \texttt{\tiny{[GM$|$GM]}}
                \texttt{L moved to (4, 4)} \\
            }
        }
    }
    & & \\ \\

    \theutterance \stepcounter{utterance}  
    & & \multicolumn{4}{p{0.6\linewidth}}{
        \cellcolor[rgb]{0.9,0.9,0.9}{
            \makecell[{{p{\linewidth}}}]{
                \texttt{\tiny{[GM$\rangle$P2]}}
                \texttt{Je bericht is doorgestuurd naar de andere speler.} \\
\\ 
\\ 
\texttt{Je speelt momenteel ronde 2 van maximaal 12.} \\
\texttt{Jullie hebben gezamenlijk 0 van 8 straffen verzameld.} \\
\texttt{De andere speler heeft een voorwerp op zijn rooster verplaatst.} \\
\\ 
\\ 
\texttt{Wat is je volgende opdracht?} \\
            }
        }
    }
    & \\ \\

    \theutterance \stepcounter{utterance}  
    & & & \multicolumn{4}{p{0.6\linewidth}}{
        \cellcolor[rgb]{0.9,0.9,0.9}{
            \makecell[{{p{\linewidth}}}]{
                \texttt{\tiny{[GM$\langle$P2]}}
                \texttt{Om mijn volgende zet te kunnen bepalen, moet ik eerst weten wat de andere speler heeft gedaan en wat de huidige status van mijn bord is (hoewel mijn objecten niet zijn verplaatst door de andere speler, is het belangrijk om de voortgang te volgen).} \\
\\ 
\texttt{Aangezien ik de actie van de andere speler niet kan zien, moet ik eerst informeren naar de status van hun bord en hun beweging.} \\
\\ 
\texttt{ZEG: Wat heb je verplaatst en op welke coördinaten staat jouw object nu? En wat zijn de huidige posities van jouw C, L en P?} \\
            }
        }
    }
     \\ \\

    \theutterance \stepcounter{utterance}  
    & & & \multicolumn{2}{p{0.3\linewidth}}{
        \cellcolor[rgb]{0.9,0.9,0.9}{
            \makecell[{{p{\linewidth}}}]{
                \texttt{\tiny{[GM$|$GM]}}
                \texttt{Om mijn volgende zet te kunnen bepalen, moet ik eerst weten wat de andere speler heeft gedaan en wat de huidige status van mijn bord is (hoewel mijn objecten niet zijn verplaatst door de andere speler, is het belangrijk om de voortgang te volgen).} \\
\\ 
\texttt{Aangezien ik de actie van de andere speler niet kan zien, moet ik eerst informeren naar de status van hun bord en hun beweging.} \\
\\ 
\texttt{ZEG: Wat heb je verplaatst en op welke coördinaten staat jouw object nu? En wat zijn de huidige posities van jouw C, L en P?} \\
            }
        }
    }
    & & \\ \\

    \theutterance \stepcounter{utterance}  
    & & & \multicolumn{2}{p{0.3\linewidth}}{
        \cellcolor[rgb]{0.9,0.9,0.9}{
            \makecell[{{p{\linewidth}}}]{
                \texttt{\tiny{[GM$|$GM]}}
                \texttt{Invalid format: head is not empty: 'Om mijn volgende zet te kunnen bepalen, moet ik eerst weten wat de andere speler heeft gedaan en wat de huidige status van mijn bord is (hoewel mijn objecten niet zijn verplaatst door de andere speler, is het belangrijk om de voortgang te volgen).} \\
\\ 
\texttt{Aangezien ik de actie van de andere speler niet kan zien, moet ik eerst informeren naar de status van hun bord en hun beweging.} \\
\\ 
\texttt{'} \\
            }
        }
    }
    & & \\ \\

    \theutterance \stepcounter{utterance}  
    & & \multicolumn{4}{p{0.6\linewidth}}{
        \cellcolor[rgb]{0.9,0.9,0.9}{
            \makecell[{{p{\linewidth}}}]{
                \texttt{\tiny{[GM$\rangle$P2]}}
                \texttt{Straf: Je bericht mag niets vóór het commando bevatten!} \\
\\ 
\texttt{Zorg ervoor dat je antwoord alleen `ZEG: <BERICHT>` of `BEWEEG: <OBJECT>, (<X>, <Y>)` bevat, en niets anders!} \\
\\ 
\texttt{Jullie hebben gezamenlijk 1 van 8 straffen verzameld.} \\
\texttt{Probeer het opnieuw!} \\
            }
        }
    }
    & \\ \\

    \theutterance \stepcounter{utterance}  
    & & & \multicolumn{4}{p{0.6\linewidth}}{
        \cellcolor[rgb]{0.9,0.9,0.9}{
            \makecell[{{p{\linewidth}}}]{
                \texttt{\tiny{[GM$\langle$P2]}}
                \texttt{ZEG: Wat heb je verplaatst en op welke coördinaten staan jouw C, L en P nu?} \\
            }
        }
    }
     \\ \\

    \theutterance \stepcounter{utterance}  
    & & & \multicolumn{2}{p{0.3\linewidth}}{
        \cellcolor[rgb]{0.9,0.9,0.9}{
            \makecell[{{p{\linewidth}}}]{
                \texttt{\tiny{[GM$|$GM]}}
                \texttt{ZEG: Wat heb je verplaatst en op welke coördinaten staan jouw C, L en P nu?} \\
            }
        }
    }
    & & \\ \\

    \theutterance \stepcounter{utterance}  
    & & \multicolumn{4}{p{0.6\linewidth}}{
        \cellcolor[rgb]{0.9,0.9,0.9}{
            \makecell[{{p{\linewidth}}}]{
                \texttt{\tiny{[P1$\langle$GM]}}
                \texttt{Verplaatst 'L' naar (4,4) met succes. Uw bijgewerkte raster ziet er als volgt uit:} \\
\texttt{```} \\
\texttt{ 1234567} \\
\texttt{╔═╤══╤══╗} \\
\texttt{║◌│◌◌│◌◌║ 1} \\
\texttt{║◌└┬─┴──╢ 2} \\
\texttt{║◌◌│◌◌◌◌║ 3} \\
\texttt{╟─┬┘L◌┌─╢ 4} \\
\texttt{║◌│◌◌◌│C║ 5} \\
\texttt{╟─┼─┬─┤P║ 6} \\
\texttt{║◌│◌│◌│◌║ 7} \\
\texttt{╚═╧═╧═╧═╝} \\
\texttt{```} \\
\texttt{Je speelt momenteel ronde 2 van maximaal 12.} \\
\texttt{Jullie hebben gezamenlijk 1 van 8 straffen verzameld.} \\
\texttt{De andere speler stuurde dit bericht:} \\
\\ 
\\ 
\texttt{"Wat heb je verplaatst en op welke coördinaten staan jouw C, L en P nu?"} \\
\\ 
\\ 
\texttt{Wat is je volgende opdracht?} \\
            }
        }
    }
    & \\ \\

    \theutterance \stepcounter{utterance}  
    & \multicolumn{4}{p{0.6\linewidth}}{
        \cellcolor[rgb]{0.9,0.9,0.9}{
            \makecell[{{p{\linewidth}}}]{
                \texttt{\tiny{[P1$\rangle$GM]}}
                \texttt{ZEG: Ik heb L verplaatst naar (4, 4). Op mijn bord staan de objecten nu op: L op (4, 4), C op (7, 5) en P op (7, 6).} \\
            }
        }
    }
    & & \\ \\

    \theutterance \stepcounter{utterance}  
    & & & \multicolumn{2}{p{0.3\linewidth}}{
        \cellcolor[rgb]{0.9,0.9,0.9}{
            \makecell[{{p{\linewidth}}}]{
                \texttt{\tiny{[GM$|$GM]}}
                \texttt{ZEG: Ik heb L verplaatst naar (4, 4). Op mijn bord staan de objecten nu op: L op (4, 4), C op (7, 5) en P op (7, 6).} \\
            }
        }
    }
    & & \\ \\

    \theutterance \stepcounter{utterance}  
    & & \multicolumn{4}{p{0.6\linewidth}}{
        \cellcolor[rgb]{0.9,0.9,0.9}{
            \makecell[{{p{\linewidth}}}]{
                \texttt{\tiny{[GM$\rangle$P2]}}
                \texttt{Je bericht is doorgestuurd naar de andere speler.} \\
\\ 
\\ 
\texttt{Je speelt momenteel ronde 3 van maximaal 12.} \\
\texttt{Jullie hebben gezamenlijk 1 van 8 straffen verzameld.} \\
\texttt{De andere speler stuurde dit bericht:} \\
\\ 
\\ 
\texttt{"Ik heb L verplaatst naar (4, 4). Op mijn bord staan de objecten nu op: L op (4, 4), C op (7, 5) en P op (7, 6)."} \\
\\ 
\\ 
\texttt{Wat is je volgende opdracht?} \\
            }
        }
    }
    & \\ \\

    \theutterance \stepcounter{utterance}  
    & & & \multicolumn{4}{p{0.6\linewidth}}{
        \cellcolor[rgb]{0.9,0.9,0.9}{
            \makecell[{{p{\linewidth}}}]{
                \texttt{\tiny{[GM$\langle$P2]}}
                \texttt{BEWEEG: L, (4, 4)} \\
            }
        }
    }
     \\ \\

    \theutterance \stepcounter{utterance}  
    & & & \multicolumn{2}{p{0.3\linewidth}}{
        \cellcolor[rgb]{0.9,0.9,0.9}{
            \makecell[{{p{\linewidth}}}]{
                \texttt{\tiny{[GM$|$GM]}}
                \texttt{BEWEEG: L, (4, 4)} \\
            }
        }
    }
    & & \\ \\

    \theutterance \stepcounter{utterance}  
    & & & \multicolumn{2}{p{0.3\linewidth}}{
        \cellcolor[rgb]{0.9,0.9,0.9}{
            \makecell[{{p{\linewidth}}}]{
                \texttt{\tiny{[GM$|$GM]}}
                \texttt{L moved to (4, 4)} \\
            }
        }
    }
    & & \\ \\

    \theutterance \stepcounter{utterance}  
    & & \multicolumn{4}{p{0.6\linewidth}}{
        \cellcolor[rgb]{0.9,0.9,0.9}{
            \makecell[{{p{\linewidth}}}]{
                \texttt{\tiny{[P1$\langle$GM]}}
                \texttt{Je bericht is doorgestuurd naar de andere speler.} \\
\\ 
\\ 
\texttt{Je speelt momenteel ronde 3 van maximaal 12.} \\
\texttt{Jullie hebben gezamenlijk 1 van 8 straffen verzameld.} \\
\texttt{De andere speler heeft een voorwerp op zijn rooster verplaatst.} \\
\\ 
\\ 
\texttt{Wat is je volgende opdracht?} \\
            }
        }
    }
    & \\ \\

    \theutterance \stepcounter{utterance}  
    & \multicolumn{4}{p{0.6\linewidth}}{
        \cellcolor[rgb]{0.9,0.9,0.9}{
            \makecell[{{p{\linewidth}}}]{
                \texttt{\tiny{[P1$\rangle$GM]}}
                \texttt{ZEG: Wat heb je verplaatst en wat zijn de nieuwe posities van jouw objecten? Ik ga daarna met mijn volgende object aan de slag.} \\
            }
        }
    }
    & & \\ \\

    \theutterance \stepcounter{utterance}  
    & & & \multicolumn{2}{p{0.3\linewidth}}{
        \cellcolor[rgb]{0.9,0.9,0.9}{
            \makecell[{{p{\linewidth}}}]{
                \texttt{\tiny{[GM$|$GM]}}
                \texttt{ZEG: Wat heb je verplaatst en wat zijn de nieuwe posities van jouw objecten? Ik ga daarna met mijn volgende object aan de slag.} \\
            }
        }
    }
    & & \\ \\

    \theutterance \stepcounter{utterance}  
    & & \multicolumn{4}{p{0.6\linewidth}}{
        \cellcolor[rgb]{0.9,0.9,0.9}{
            \makecell[{{p{\linewidth}}}]{
                \texttt{\tiny{[GM$\rangle$P2]}}
                \texttt{Verplaatst 'L' naar (4,4) met succes. Uw bijgewerkte raster ziet er als volgt uit:} \\
\texttt{```} \\
\texttt{ 1234567} \\
\texttt{╔═╤══╤══╗} \\
\texttt{║◌│◌◌│◌◌║ 1} \\
\texttt{║◌└┬─┴──╢ 2} \\
\texttt{║◌C│◌◌◌◌║ 3} \\
\texttt{╟─┬┘L◌┌─╢ 4} \\
\texttt{║◌│◌◌◌│◌║ 5} \\
\texttt{╟─┼─┬─┤◌║ 6} \\
\texttt{║◌│P│◌│◌║ 7} \\
\texttt{╚═╧═╧═╧═╝} \\
\texttt{```} \\
\texttt{Je speelt momenteel ronde 4 van maximaal 12.} \\
\texttt{Jullie hebben gezamenlijk 1 van 8 straffen verzameld.} \\
\texttt{De andere speler stuurde dit bericht:} \\
\\ 
\\ 
\texttt{"Wat heb je verplaatst en wat zijn de nieuwe posities van jouw objecten? Ik ga daarna met mijn volgende object aan de slag."} \\
\\ 
\\ 
\texttt{Wat is je volgende opdracht?} \\
            }
        }
    }
    & \\ \\

    \theutterance \stepcounter{utterance}  
    & & & \multicolumn{4}{p{0.6\linewidth}}{
        \cellcolor[rgb]{0.9,0.9,0.9}{
            \makecell[{{p{\linewidth}}}]{
                \texttt{\tiny{[GM$\langle$P2]}}
                \texttt{ZEG: Ik heb L verplaatst naar (4, 4). Mijn objecten staan nu op: L op (4, 4), C op (2, 3) en P op (3, 7).} \\
            }
        }
    }
     \\ \\

    \theutterance \stepcounter{utterance}  
    & & & \multicolumn{2}{p{0.3\linewidth}}{
        \cellcolor[rgb]{0.9,0.9,0.9}{
            \makecell[{{p{\linewidth}}}]{
                \texttt{\tiny{[GM$|$GM]}}
                \texttt{ZEG: Ik heb L verplaatst naar (4, 4). Mijn objecten staan nu op: L op (4, 4), C op (2, 3) en P op (3, 7).} \\
            }
        }
    }
    & & \\ \\

    \theutterance \stepcounter{utterance}  
    & & \multicolumn{4}{p{0.6\linewidth}}{
        \cellcolor[rgb]{0.9,0.9,0.9}{
            \makecell[{{p{\linewidth}}}]{
                \texttt{\tiny{[P1$\langle$GM]}}
                \texttt{Je bericht is doorgestuurd naar de andere speler.} \\
\\ 
\\ 
\texttt{Je speelt momenteel ronde 4 van maximaal 12.} \\
\texttt{Jullie hebben gezamenlijk 1 van 8 straffen verzameld.} \\
\texttt{De andere speler stuurde dit bericht:} \\
\\ 
\\ 
\texttt{"Ik heb L verplaatst naar (4, 4). Mijn objecten staan nu op: L op (4, 4), C op (2, 3) en P op (3, 7)."} \\
\\ 
\\ 
\texttt{Wat is je volgende opdracht?} \\
            }
        }
    }
    & \\ \\

    \theutterance \stepcounter{utterance}  
    & \multicolumn{4}{p{0.6\linewidth}}{
        \cellcolor[rgb]{0.9,0.9,0.9}{
            \makecell[{{p{\linewidth}}}]{
                \texttt{\tiny{[P1$\rangle$GM]}}
                \texttt{BEWEEG: C, (4, 4)} \\
            }
        }
    }
    & & \\ \\

    \theutterance \stepcounter{utterance}  
    & & & \multicolumn{2}{p{0.3\linewidth}}{
        \cellcolor[rgb]{0.9,0.9,0.9}{
            \makecell[{{p{\linewidth}}}]{
                \texttt{\tiny{[GM$|$GM]}}
                \texttt{BEWEEG: C, (4, 4)} \\
            }
        }
    }
    & & \\ \\

    \theutterance \stepcounter{utterance}  
    & & & \multicolumn{2}{p{0.3\linewidth}}{
        \cellcolor[rgb]{0.9,0.9,0.9}{
            \makecell[{{p{\linewidth}}}]{
                \texttt{\tiny{[GM$|$GM]}}
                \texttt{Straf: (4,4) is niet leeg, maar bevat 'L'.} \\
\\ 
\texttt{Jullie hebben gezamenlijk 2 van 8 straffen verzameld.} \\
\texttt{Probeer het opnieuw!} \\
            }
        }
    }
    & & \\ \\

    \theutterance \stepcounter{utterance}  
    & & \multicolumn{4}{p{0.6\linewidth}}{
        \cellcolor[rgb]{0.9,0.9,0.9}{
            \makecell[{{p{\linewidth}}}]{
                \texttt{\tiny{[P1$\langle$GM]}}
                \texttt{Straf: (4,4) is niet leeg, maar bevat 'L'.} \\
\\ 
\texttt{Jullie hebben gezamenlijk 2 van 8 straffen verzameld.} \\
\texttt{Probeer het opnieuw!} \\
            }
        }
    }
    & \\ \\

    \theutterance \stepcounter{utterance}  
    & \multicolumn{4}{p{0.6\linewidth}}{
        \cellcolor[rgb]{0.9,0.9,0.9}{
            \makecell[{{p{\linewidth}}}]{
                \texttt{\tiny{[P1$\rangle$GM]}}
                \texttt{BEWEEG: C, (4, 5)} \\
            }
        }
    }
    & & \\ \\

    \theutterance \stepcounter{utterance}  
    & & & \multicolumn{2}{p{0.3\linewidth}}{
        \cellcolor[rgb]{0.9,0.9,0.9}{
            \makecell[{{p{\linewidth}}}]{
                \texttt{\tiny{[GM$|$GM]}}
                \texttt{BEWEEG: C, (4, 5)} \\
            }
        }
    }
    & & \\ \\

    \theutterance \stepcounter{utterance}  
    & & & \multicolumn{2}{p{0.3\linewidth}}{
        \cellcolor[rgb]{0.9,0.9,0.9}{
            \makecell[{{p{\linewidth}}}]{
                \texttt{\tiny{[GM$|$GM]}}
                \texttt{C moved to (4, 5)} \\
            }
        }
    }
    & & \\ \\

    \theutterance \stepcounter{utterance}  
    & & \multicolumn{4}{p{0.6\linewidth}}{
        \cellcolor[rgb]{0.9,0.9,0.9}{
            \makecell[{{p{\linewidth}}}]{
                \texttt{\tiny{[GM$\rangle$P2]}}
                \texttt{Je bericht is doorgestuurd naar de andere speler.} \\
\\ 
\\ 
\texttt{Je speelt momenteel ronde 5 van maximaal 12.} \\
\texttt{Jullie hebben gezamenlijk 2 van 8 straffen verzameld.} \\
\texttt{De andere speler heeft een voorwerp op zijn rooster verplaatst.} \\
\\ 
\\ 
\texttt{Wat is je volgende opdracht?} \\
            }
        }
    }
    & \\ \\

    \theutterance \stepcounter{utterance}  
    & & & \multicolumn{4}{p{0.6\linewidth}}{
        \cellcolor[rgb]{0.9,0.9,0.9}{
            \makecell[{{p{\linewidth}}}]{
                \texttt{\tiny{[GM$\langle$P2]}}
                \texttt{ZEG: Wat heb je verplaatst en wat zijn de nieuwe posities van jouw C en P?} \\
            }
        }
    }
     \\ \\

    \theutterance \stepcounter{utterance}  
    & & & \multicolumn{2}{p{0.3\linewidth}}{
        \cellcolor[rgb]{0.9,0.9,0.9}{
            \makecell[{{p{\linewidth}}}]{
                \texttt{\tiny{[GM$|$GM]}}
                \texttt{ZEG: Wat heb je verplaatst en wat zijn de nieuwe posities van jouw C en P?} \\
            }
        }
    }
    & & \\ \\

    \theutterance \stepcounter{utterance}  
    & & \multicolumn{4}{p{0.6\linewidth}}{
        \cellcolor[rgb]{0.9,0.9,0.9}{
            \makecell[{{p{\linewidth}}}]{
                \texttt{\tiny{[P1$\langle$GM]}}
                \texttt{Verplaatst 'C' naar (4,5) met succes. Uw bijgewerkte raster ziet er als volgt uit:} \\
\texttt{```} \\
\texttt{ 1234567} \\
\texttt{╔═╤══╤══╗} \\
\texttt{║◌│◌◌│◌◌║ 1} \\
\texttt{║◌└┬─┴──╢ 2} \\
\texttt{║◌◌│◌◌◌◌║ 3} \\
\texttt{╟─┬┘L◌┌─╢ 4} \\
\texttt{║◌│◌C◌│◌║ 5} \\
\texttt{╟─┼─┬─┤P║ 6} \\
\texttt{║◌│◌│◌│◌║ 7} \\
\texttt{╚═╧═╧═╧═╝} \\
\texttt{```} \\
\texttt{Je speelt momenteel ronde 5 van maximaal 12.} \\
\texttt{Jullie hebben gezamenlijk 2 van 8 straffen verzameld.} \\
\texttt{De andere speler stuurde dit bericht:} \\
\\ 
\\ 
\texttt{"Wat heb je verplaatst en wat zijn de nieuwe posities van jouw C en P?"} \\
\\ 
\\ 
\texttt{Wat is je volgende opdracht?} \\
            }
        }
    }
    & \\ \\

    \theutterance \stepcounter{utterance}  
    & \multicolumn{4}{p{0.6\linewidth}}{
        \cellcolor[rgb]{0.9,0.9,0.9}{
            \makecell[{{p{\linewidth}}}]{
                \texttt{\tiny{[P1$\rangle$GM]}}
                \texttt{ZEG: Ik heb C verplaatst naar (4, 5). Mijn objecten staan nu op: L op (4, 4), C op (4, 5) en P op (7, 6).} \\
            }
        }
    }
    & & \\ \\

    \theutterance \stepcounter{utterance}  
    & & & \multicolumn{2}{p{0.3\linewidth}}{
        \cellcolor[rgb]{0.9,0.9,0.9}{
            \makecell[{{p{\linewidth}}}]{
                \texttt{\tiny{[GM$|$GM]}}
                \texttt{ZEG: Ik heb C verplaatst naar (4, 5). Mijn objecten staan nu op: L op (4, 4), C op (4, 5) en P op (7, 6).} \\
            }
        }
    }
    & & \\ \\

    \theutterance \stepcounter{utterance}  
    & & \multicolumn{4}{p{0.6\linewidth}}{
        \cellcolor[rgb]{0.9,0.9,0.9}{
            \makecell[{{p{\linewidth}}}]{
                \texttt{\tiny{[GM$\rangle$P2]}}
                \texttt{Je bericht is doorgestuurd naar de andere speler.} \\
\\ 
\\ 
\texttt{Je speelt momenteel ronde 6 van maximaal 12.} \\
\texttt{Jullie hebben gezamenlijk 2 van 8 straffen verzameld.} \\
\texttt{De andere speler stuurde dit bericht:} \\
\\ 
\\ 
\texttt{"Ik heb C verplaatst naar (4, 5). Mijn objecten staan nu op: L op (4, 4), C op (4, 5) en P op (7, 6)."} \\
\\ 
\\ 
\texttt{Wat is je volgende opdracht?} \\
            }
        }
    }
    & \\ \\

    \theutterance \stepcounter{utterance}  
    & & & \multicolumn{4}{p{0.6\linewidth}}{
        \cellcolor[rgb]{0.9,0.9,0.9}{
            \makecell[{{p{\linewidth}}}]{
                \texttt{\tiny{[GM$\langle$P2]}}
                \texttt{BEWEEG: C, (4, 4)} \\
            }
        }
    }
     \\ \\

    \theutterance \stepcounter{utterance}  
    & & & \multicolumn{2}{p{0.3\linewidth}}{
        \cellcolor[rgb]{0.9,0.9,0.9}{
            \makecell[{{p{\linewidth}}}]{
                \texttt{\tiny{[GM$|$GM]}}
                \texttt{BEWEEG: C, (4, 4)} \\
            }
        }
    }
    & & \\ \\

    \theutterance \stepcounter{utterance}  
    & & & \multicolumn{2}{p{0.3\linewidth}}{
        \cellcolor[rgb]{0.9,0.9,0.9}{
            \makecell[{{p{\linewidth}}}]{
                \texttt{\tiny{[GM$|$GM]}}
                \texttt{Straf: (4,4) is niet leeg, maar bevat 'L'.} \\
\\ 
\texttt{Jullie hebben gezamenlijk 3 van 8 straffen verzameld.} \\
\texttt{Probeer het opnieuw!} \\
            }
        }
    }
    & & \\ \\

    \theutterance \stepcounter{utterance}  
    & & \multicolumn{4}{p{0.6\linewidth}}{
        \cellcolor[rgb]{0.9,0.9,0.9}{
            \makecell[{{p{\linewidth}}}]{
                \texttt{\tiny{[GM$\rangle$P2]}}
                \texttt{Straf: (4,4) is niet leeg, maar bevat 'L'.} \\
\\ 
\texttt{Jullie hebben gezamenlijk 3 van 8 straffen verzameld.} \\
\texttt{Probeer het opnieuw!} \\
            }
        }
    }
    & \\ \\

    \theutterance \stepcounter{utterance}  
    & & & \multicolumn{4}{p{0.6\linewidth}}{
        \cellcolor[rgb]{0.9,0.9,0.9}{
            \makecell[{{p{\linewidth}}}]{
                \texttt{\tiny{[GM$\langle$P2]}}
                \texttt{BEWEEG: C, (4, 5)} \\
            }
        }
    }
     \\ \\

    \theutterance \stepcounter{utterance}  
    & & & \multicolumn{2}{p{0.3\linewidth}}{
        \cellcolor[rgb]{0.9,0.9,0.9}{
            \makecell[{{p{\linewidth}}}]{
                \texttt{\tiny{[GM$|$GM]}}
                \texttt{BEWEEG: C, (4, 5)} \\
            }
        }
    }
    & & \\ \\

    \theutterance \stepcounter{utterance}  
    & & & \multicolumn{2}{p{0.3\linewidth}}{
        \cellcolor[rgb]{0.9,0.9,0.9}{
            \makecell[{{p{\linewidth}}}]{
                \texttt{\tiny{[GM$|$GM]}}
                \texttt{C moved to (4, 5)} \\
            }
        }
    }
    & & \\ \\

    \theutterance \stepcounter{utterance}  
    & & \multicolumn{4}{p{0.6\linewidth}}{
        \cellcolor[rgb]{0.9,0.9,0.9}{
            \makecell[{{p{\linewidth}}}]{
                \texttt{\tiny{[P1$\langle$GM]}}
                \texttt{Je bericht is doorgestuurd naar de andere speler.} \\
\\ 
\\ 
\texttt{Je speelt momenteel ronde 6 van maximaal 12.} \\
\texttt{Jullie hebben gezamenlijk 3 van 8 straffen verzameld.} \\
\texttt{De andere speler heeft een voorwerp op zijn rooster verplaatst.} \\
\\ 
\\ 
\texttt{Wat is je volgende opdracht?} \\
            }
        }
    }
    & \\ \\

    \theutterance \stepcounter{utterance}  
    & \multicolumn{4}{p{0.6\linewidth}}{
        \cellcolor[rgb]{0.9,0.9,0.9}{
            \makecell[{{p{\linewidth}}}]{
                \texttt{\tiny{[P1$\rangle$GM]}}
                \texttt{ZEG: Wat heb je verplaatst en wat zijn de nieuwe posities van jouw C en P? Ik ga nu proberen om mijn P naar (4, 4) te verplaatsen.} \\
            }
        }
    }
    & & \\ \\

    \theutterance \stepcounter{utterance}  
    & & & \multicolumn{2}{p{0.3\linewidth}}{
        \cellcolor[rgb]{0.9,0.9,0.9}{
            \makecell[{{p{\linewidth}}}]{
                \texttt{\tiny{[GM$|$GM]}}
                \texttt{ZEG: Wat heb je verplaatst en wat zijn de nieuwe posities van jouw C en P? Ik ga nu proberen om mijn P naar (4, 4) te verplaatsen.} \\
            }
        }
    }
    & & \\ \\

    \theutterance \stepcounter{utterance}  
    & & \multicolumn{4}{p{0.6\linewidth}}{
        \cellcolor[rgb]{0.9,0.9,0.9}{
            \makecell[{{p{\linewidth}}}]{
                \texttt{\tiny{[GM$\rangle$P2]}}
                \texttt{Verplaatst 'C' naar (4,5) met succes. Uw bijgewerkte raster ziet er als volgt uit:} \\
\texttt{```} \\
\texttt{ 1234567} \\
\texttt{╔═╤══╤══╗} \\
\texttt{║◌│◌◌│◌◌║ 1} \\
\texttt{║◌└┬─┴──╢ 2} \\
\texttt{║◌◌│◌◌◌◌║ 3} \\
\texttt{╟─┬┘L◌┌─╢ 4} \\
\texttt{║◌│◌C◌│◌║ 5} \\
\texttt{╟─┼─┬─┤◌║ 6} \\
\texttt{║◌│P│◌│◌║ 7} \\
\texttt{╚═╧═╧═╧═╝} \\
\texttt{```} \\
\texttt{Je speelt momenteel ronde 7 van maximaal 12.} \\
\texttt{Jullie hebben gezamenlijk 3 van 8 straffen verzameld.} \\
\texttt{De andere speler stuurde dit bericht:} \\
\\ 
\\ 
\texttt{"Wat heb je verplaatst en wat zijn de nieuwe posities van jouw C en P? Ik ga nu proberen om mijn P naar (4, 4) te verplaatsen."} \\
\\ 
\\ 
\texttt{Wat is je volgende opdracht?} \\
            }
        }
    }
    & \\ \\

    \theutterance \stepcounter{utterance}  
    & & & \multicolumn{4}{p{0.6\linewidth}}{
        \cellcolor[rgb]{0.9,0.9,0.9}{
            \makecell[{{p{\linewidth}}}]{
                \texttt{\tiny{[GM$\langle$P2]}}
                \texttt{ZEG: Ik heb C verplaatst naar (4, 5). Mijn objecten staan nu op: L op (4, 4), C op (4, 5) en P op (3, 7).} \\
            }
        }
    }
     \\ \\

    \theutterance \stepcounter{utterance}  
    & & & \multicolumn{2}{p{0.3\linewidth}}{
        \cellcolor[rgb]{0.9,0.9,0.9}{
            \makecell[{{p{\linewidth}}}]{
                \texttt{\tiny{[GM$|$GM]}}
                \texttt{ZEG: Ik heb C verplaatst naar (4, 5). Mijn objecten staan nu op: L op (4, 4), C op (4, 5) en P op (3, 7).} \\
            }
        }
    }
    & & \\ \\

    \theutterance \stepcounter{utterance}  
    & & \multicolumn{4}{p{0.6\linewidth}}{
        \cellcolor[rgb]{0.9,0.9,0.9}{
            \makecell[{{p{\linewidth}}}]{
                \texttt{\tiny{[P1$\langle$GM]}}
                \texttt{Je bericht is doorgestuurd naar de andere speler.} \\
\\ 
\\ 
\texttt{Je speelt momenteel ronde 7 van maximaal 12.} \\
\texttt{Jullie hebben gezamenlijk 3 van 8 straffen verzameld.} \\
\texttt{De andere speler stuurde dit bericht:} \\
\\ 
\\ 
\texttt{"Ik heb C verplaatst naar (4, 5). Mijn objecten staan nu op: L op (4, 4), C op (4, 5) en P op (3, 7)."} \\
\\ 
\\ 
\texttt{Wat is je volgende opdracht?} \\
            }
        }
    }
    & \\ \\

    \theutterance \stepcounter{utterance}  
    & \multicolumn{4}{p{0.6\linewidth}}{
        \cellcolor[rgb]{0.9,0.9,0.9}{
            \makecell[{{p{\linewidth}}}]{
                \texttt{\tiny{[P1$\rangle$GM]}}
                \texttt{BEWEEG: P, (4, 4)} \\
            }
        }
    }
    & & \\ \\

    \theutterance \stepcounter{utterance}  
    & & & \multicolumn{2}{p{0.3\linewidth}}{
        \cellcolor[rgb]{0.9,0.9,0.9}{
            \makecell[{{p{\linewidth}}}]{
                \texttt{\tiny{[GM$|$GM]}}
                \texttt{BEWEEG: P, (4, 4)} \\
            }
        }
    }
    & & \\ \\

    \theutterance \stepcounter{utterance}  
    & & & \multicolumn{2}{p{0.3\linewidth}}{
        \cellcolor[rgb]{0.9,0.9,0.9}{
            \makecell[{{p{\linewidth}}}]{
                \texttt{\tiny{[GM$|$GM]}}
                \texttt{Straf: (4,4) is niet leeg, maar bevat 'L'.} \\
\\ 
\texttt{Jullie hebben gezamenlijk 4 van 8 straffen verzameld.} \\
\texttt{Probeer het opnieuw!} \\
            }
        }
    }
    & & \\ \\

    \theutterance \stepcounter{utterance}  
    & & \multicolumn{4}{p{0.6\linewidth}}{
        \cellcolor[rgb]{0.9,0.9,0.9}{
            \makecell[{{p{\linewidth}}}]{
                \texttt{\tiny{[P1$\langle$GM]}}
                \texttt{Straf: (4,4) is niet leeg, maar bevat 'L'.} \\
\\ 
\texttt{Jullie hebben gezamenlijk 4 van 8 straffen verzameld.} \\
\texttt{Probeer het opnieuw!} \\
            }
        }
    }
    & \\ \\

    \theutterance \stepcounter{utterance}  
    & \multicolumn{4}{p{0.6\linewidth}}{
        \cellcolor[rgb]{0.9,0.9,0.9}{
            \makecell[{{p{\linewidth}}}]{
                \texttt{\tiny{[P1$\rangle$GM]}}
                \texttt{BEWEEG: P, (4, 6)} \\
            }
        }
    }
    & & \\ \\

    \theutterance \stepcounter{utterance}  
    & & & \multicolumn{2}{p{0.3\linewidth}}{
        \cellcolor[rgb]{0.9,0.9,0.9}{
            \makecell[{{p{\linewidth}}}]{
                \texttt{\tiny{[GM$|$GM]}}
                \texttt{BEWEEG: P, (4, 6)} \\
            }
        }
    }
    & & \\ \\

    \theutterance \stepcounter{utterance}  
    & & & \multicolumn{2}{p{0.3\linewidth}}{
        \cellcolor[rgb]{0.9,0.9,0.9}{
            \makecell[{{p{\linewidth}}}]{
                \texttt{\tiny{[GM$|$GM]}}
                \texttt{Straf: (4,6) is niet leeg, maar bevat '┬'.} \\
\\ 
\texttt{Jullie hebben gezamenlijk 5 van 8 straffen verzameld.} \\
\texttt{Probeer het opnieuw!} \\
            }
        }
    }
    & & \\ \\

    \theutterance \stepcounter{utterance}  
    & & \multicolumn{4}{p{0.6\linewidth}}{
        \cellcolor[rgb]{0.9,0.9,0.9}{
            \makecell[{{p{\linewidth}}}]{
                \texttt{\tiny{[P1$\langle$GM]}}
                \texttt{Straf: (4,6) is niet leeg, maar bevat '┬'.} \\
\\ 
\texttt{Jullie hebben gezamenlijk 5 van 8 straffen verzameld.} \\
\texttt{Probeer het opnieuw!} \\
            }
        }
    }
    & \\ \\

    \theutterance \stepcounter{utterance}  
    & \multicolumn{4}{p{0.6\linewidth}}{
        \cellcolor[rgb]{0.9,0.9,0.9}{
            \makecell[{{p{\linewidth}}}]{
                \texttt{\tiny{[P1$\rangle$GM]}}
                \texttt{BEWEEG: P, (4, 7)} \\
            }
        }
    }
    & & \\ \\

    \theutterance \stepcounter{utterance}  
    & & & \multicolumn{2}{p{0.3\linewidth}}{
        \cellcolor[rgb]{0.9,0.9,0.9}{
            \makecell[{{p{\linewidth}}}]{
                \texttt{\tiny{[GM$|$GM]}}
                \texttt{BEWEEG: P, (4, 7)} \\
            }
        }
    }
    & & \\ \\

    \theutterance \stepcounter{utterance}  
    & & & \multicolumn{2}{p{0.3\linewidth}}{
        \cellcolor[rgb]{0.9,0.9,0.9}{
            \makecell[{{p{\linewidth}}}]{
                \texttt{\tiny{[GM$|$GM]}}
                \texttt{Straf: (4,7) is niet leeg, maar bevat '│'.} \\
\\ 
\texttt{Jullie hebben gezamenlijk 6 van 8 straffen verzameld.} \\
\texttt{Probeer het opnieuw!} \\
            }
        }
    }
    & & \\ \\

    \theutterance \stepcounter{utterance}  
    & & \multicolumn{4}{p{0.6\linewidth}}{
        \cellcolor[rgb]{0.9,0.9,0.9}{
            \makecell[{{p{\linewidth}}}]{
                \texttt{\tiny{[P1$\langle$GM]}}
                \texttt{Straf: (4,7) is niet leeg, maar bevat '│'.} \\
\\ 
\texttt{Jullie hebben gezamenlijk 6 van 8 straffen verzameld.} \\
\texttt{Probeer het opnieuw!} \\
            }
        }
    }
    & \\ \\

    \theutterance \stepcounter{utterance}  
    & \multicolumn{4}{p{0.6\linewidth}}{
        \cellcolor[rgb]{0.9,0.9,0.9}{
            \makecell[{{p{\linewidth}}}]{
                \texttt{\tiny{[P1$\rangle$GM]}}
                \texttt{BEWEEG: P, (5, 7)} \\
            }
        }
    }
    & & \\ \\

    \theutterance \stepcounter{utterance}  
    & & & \multicolumn{2}{p{0.3\linewidth}}{
        \cellcolor[rgb]{0.9,0.9,0.9}{
            \makecell[{{p{\linewidth}}}]{
                \texttt{\tiny{[GM$|$GM]}}
                \texttt{BEWEEG: P, (5, 7)} \\
            }
        }
    }
    & & \\ \\

    \theutterance \stepcounter{utterance}  
    & & & \multicolumn{2}{p{0.3\linewidth}}{
        \cellcolor[rgb]{0.9,0.9,0.9}{
            \makecell[{{p{\linewidth}}}]{
                \texttt{\tiny{[GM$|$GM]}}
                \texttt{P moved to (5, 7)} \\
            }
        }
    }
    & & \\ \\

    \theutterance \stepcounter{utterance}  
    & & \multicolumn{4}{p{0.6\linewidth}}{
        \cellcolor[rgb]{0.9,0.9,0.9}{
            \makecell[{{p{\linewidth}}}]{
                \texttt{\tiny{[GM$\rangle$P2]}}
                \texttt{Je bericht is doorgestuurd naar de andere speler.} \\
\\ 
\\ 
\texttt{Je speelt momenteel ronde 8 van maximaal 12.} \\
\texttt{Jullie hebben gezamenlijk 6 van 8 straffen verzameld.} \\
\texttt{De andere speler heeft een voorwerp op zijn rooster verplaatst.} \\
\\ 
\\ 
\texttt{Wat is je volgende opdracht?} \\
            }
        }
    }
    & \\ \\

    \theutterance \stepcounter{utterance}  
    & & & \multicolumn{4}{p{0.6\linewidth}}{
        \cellcolor[rgb]{0.9,0.9,0.9}{
            \makecell[{{p{\linewidth}}}]{
                \texttt{\tiny{[GM$\langle$P2]}}
                \texttt{ZEG: Wat heb je verplaatst en wat zijn de nieuwe posities van jouw L, C en P?} \\
            }
        }
    }
     \\ \\

    \theutterance \stepcounter{utterance}  
    & & & \multicolumn{2}{p{0.3\linewidth}}{
        \cellcolor[rgb]{0.9,0.9,0.9}{
            \makecell[{{p{\linewidth}}}]{
                \texttt{\tiny{[GM$|$GM]}}
                \texttt{ZEG: Wat heb je verplaatst en wat zijn de nieuwe posities van jouw L, C en P?} \\
            }
        }
    }
    & & \\ \\

    \theutterance \stepcounter{utterance}  
    & & \multicolumn{4}{p{0.6\linewidth}}{
        \cellcolor[rgb]{0.9,0.9,0.9}{
            \makecell[{{p{\linewidth}}}]{
                \texttt{\tiny{[P1$\langle$GM]}}
                \texttt{Verplaatst 'P' naar (5,7) met succes. Uw bijgewerkte raster ziet er als volgt uit:} \\
\texttt{```} \\
\texttt{ 1234567} \\
\texttt{╔═╤══╤══╗} \\
\texttt{║◌│◌◌│◌◌║ 1} \\
\texttt{║◌└┬─┴──╢ 2} \\
\texttt{║◌◌│◌◌◌◌║ 3} \\
\texttt{╟─┬┘L◌┌─╢ 4} \\
\texttt{║◌│◌C◌│◌║ 5} \\
\texttt{╟─┼─┬─┤◌║ 6} \\
\texttt{║◌│◌│P│◌║ 7} \\
\texttt{╚═╧═╧═╧═╝} \\
\texttt{```} \\
\texttt{Je speelt momenteel ronde 8 van maximaal 12.} \\
\texttt{Jullie hebben gezamenlijk 6 van 8 straffen verzameld.} \\
\texttt{De andere speler stuurde dit bericht:} \\
\\ 
\\ 
\texttt{"Wat heb je verplaatst en wat zijn de nieuwe posities van jouw L, C en P?"} \\
\\ 
\\ 
\texttt{Wat is je volgende opdracht?} \\
            }
        }
    }
    & \\ \\

    \theutterance \stepcounter{utterance}  
    & \multicolumn{4}{p{0.6\linewidth}}{
        \cellcolor[rgb]{0.9,0.9,0.9}{
            \makecell[{{p{\linewidth}}}]{
                \texttt{\tiny{[P1$\rangle$GM]}}
                \texttt{ZEG: Ik heb P verplaatst naar (5, 7). Mijn objecten staan nu op: L op (4, 4), C op (4, 5) en P op (5, 7).} \\
            }
        }
    }
    & & \\ \\

    \theutterance \stepcounter{utterance}  
    & & & \multicolumn{2}{p{0.3\linewidth}}{
        \cellcolor[rgb]{0.9,0.9,0.9}{
            \makecell[{{p{\linewidth}}}]{
                \texttt{\tiny{[GM$|$GM]}}
                \texttt{ZEG: Ik heb P verplaatst naar (5, 7). Mijn objecten staan nu op: L op (4, 4), C op (4, 5) en P op (5, 7).} \\
            }
        }
    }
    & & \\ \\

    \theutterance \stepcounter{utterance}  
    & & \multicolumn{4}{p{0.6\linewidth}}{
        \cellcolor[rgb]{0.9,0.9,0.9}{
            \makecell[{{p{\linewidth}}}]{
                \texttt{\tiny{[GM$\rangle$P2]}}
                \texttt{Je bericht is doorgestuurd naar de andere speler.} \\
\\ 
\\ 
\texttt{Je speelt momenteel ronde 9 van maximaal 12.} \\
\texttt{Jullie hebben gezamenlijk 6 van 8 straffen verzameld.} \\
\texttt{De andere speler stuurde dit bericht:} \\
\\ 
\\ 
\texttt{"Ik heb P verplaatst naar (5, 7). Mijn objecten staan nu op: L op (4, 4), C op (4, 5) en P op (5, 7)."} \\
\\ 
\\ 
\texttt{Wat is je volgende opdracht?} \\
            }
        }
    }
    & \\ \\

    \theutterance \stepcounter{utterance}  
    & & & \multicolumn{4}{p{0.6\linewidth}}{
        \cellcolor[rgb]{0.9,0.9,0.9}{
            \makecell[{{p{\linewidth}}}]{
                \texttt{\tiny{[GM$\langle$P2]}}
                \texttt{BEWEEG: P, (4, 4)} \\
            }
        }
    }
     \\ \\

    \theutterance \stepcounter{utterance}  
    & & & \multicolumn{2}{p{0.3\linewidth}}{
        \cellcolor[rgb]{0.9,0.9,0.9}{
            \makecell[{{p{\linewidth}}}]{
                \texttt{\tiny{[GM$|$GM]}}
                \texttt{BEWEEG: P, (4, 4)} \\
            }
        }
    }
    & & \\ \\

    \theutterance \stepcounter{utterance}  
    & & & \multicolumn{2}{p{0.3\linewidth}}{
        \cellcolor[rgb]{0.9,0.9,0.9}{
            \makecell[{{p{\linewidth}}}]{
                \texttt{\tiny{[GM$|$GM]}}
                \texttt{Straf: (4,4) is niet leeg, maar bevat 'L'.} \\
\\ 
\texttt{Jullie hebben gezamenlijk 7 van 8 straffen verzameld.} \\
\texttt{Probeer het opnieuw!} \\
            }
        }
    }
    & & \\ \\

    \theutterance \stepcounter{utterance}  
    & & \multicolumn{4}{p{0.6\linewidth}}{
        \cellcolor[rgb]{0.9,0.9,0.9}{
            \makecell[{{p{\linewidth}}}]{
                \texttt{\tiny{[GM$\rangle$P2]}}
                \texttt{Straf: (4,4) is niet leeg, maar bevat 'L'.} \\
\\ 
\texttt{Jullie hebben gezamenlijk 7 van 8 straffen verzameld.} \\
\texttt{Probeer het opnieuw!} \\
            }
        }
    }
    & \\ \\

    \theutterance \stepcounter{utterance}  
    & & & \multicolumn{4}{p{0.6\linewidth}}{
        \cellcolor[rgb]{0.9,0.9,0.9}{
            \makecell[{{p{\linewidth}}}]{
                \texttt{\tiny{[GM$\langle$P2]}}
                \texttt{BEWEEG: P, (4, 5)} \\
            }
        }
    }
     \\ \\

    \theutterance \stepcounter{utterance}  
    & & & \multicolumn{2}{p{0.3\linewidth}}{
        \cellcolor[rgb]{0.9,0.9,0.9}{
            \makecell[{{p{\linewidth}}}]{
                \texttt{\tiny{[GM$|$GM]}}
                \texttt{BEWEEG: P, (4, 5)} \\
            }
        }
    }
    & & \\ \\

    \theutterance \stepcounter{utterance}  
    & & & \multicolumn{2}{p{0.3\linewidth}}{
        \cellcolor[rgb]{0.9,0.9,0.9}{
            \makecell[{{p{\linewidth}}}]{
                \texttt{\tiny{[GM$|$GM]}}
                \texttt{Straf: (4,5) is niet leeg, maar bevat 'C'.} \\
\\ 
\texttt{Jullie hebben gezamenlijk 8 van 8 straffen verzameld.} \\
\texttt{Probeer het opnieuw!} \\
            }
        }
    }
    & & \\ \\

    \theutterance \stepcounter{utterance}  
    & & \multicolumn{4}{p{0.6\linewidth}}{
        \cellcolor[rgb]{0.9,0.9,0.9}{
            \makecell[{{p{\linewidth}}}]{
                \texttt{\tiny{[GM$\rangle$P2]}}
                \texttt{Straf: (4,5) is niet leeg, maar bevat 'C'.} \\
\\ 
\texttt{Jullie hebben gezamenlijk 8 van 8 straffen verzameld.} \\
\texttt{Probeer het opnieuw!} \\
            }
        }
    }
    & \\ \\

    \theutterance \stepcounter{utterance}  
    & & & \multicolumn{4}{p{0.6\linewidth}}{
        \cellcolor[rgb]{0.9,0.9,0.9}{
            \makecell[{{p{\linewidth}}}]{
                \texttt{\tiny{[GM$\langle$P2]}}
                \texttt{BEWEEG: P, (4, 6)} \\
            }
        }
    }
     \\ \\

    \theutterance \stepcounter{utterance}  
    & & & \multicolumn{2}{p{0.3\linewidth}}{
        \cellcolor[rgb]{0.9,0.9,0.9}{
            \makecell[{{p{\linewidth}}}]{
                \texttt{\tiny{[GM$|$GM]}}
                \texttt{BEWEEG: P, (4, 6)} \\
            }
        }
    }
    & & \\ \\

    \theutterance \stepcounter{utterance}  
    & & & \multicolumn{2}{p{0.3\linewidth}}{
        \cellcolor[rgb]{0.9,0.9,0.9}{
            \makecell[{{p{\linewidth}}}]{
                \texttt{\tiny{[GM$|$GM]}}
                \texttt{Straf: (4,6) is niet leeg, maar bevat '┬'.} \\
\\ 
\texttt{Jullie hebben gezamenlijk 9 van 8 straffen verzameld.} \\
\texttt{Probeer het opnieuw!} \\
            }
        }
    }
    & & \\ \\

    \theutterance \stepcounter{utterance}  
    & & \multicolumn{4}{p{0.6\linewidth}}{
        \cellcolor[rgb]{0.9,0.9,0.9}{
            \makecell[{{p{\linewidth}}}]{
                \texttt{\tiny{[GM$\rangle$P2]}}
                \texttt{Straf: (4,6) is niet leeg, maar bevat '┬'.} \\
\\ 
\texttt{Jullie hebben gezamenlijk 9 van 8 straffen verzameld.} \\
\texttt{Probeer het opnieuw!} \\
            }
        }
    }
    & \\ \\

    \theutterance \stepcounter{utterance}  
    & & & \multicolumn{4}{p{0.6\linewidth}}{
        \cellcolor[rgb]{0.9,0.9,0.9}{
            \makecell[{{p{\linewidth}}}]{
                \texttt{\tiny{[GM$\langle$P2]}}
                \texttt{BEWEEG: P, (3, 6)} \\
            }
        }
    }
     \\ \\

    \theutterance \stepcounter{utterance}  
    & & & \multicolumn{2}{p{0.3\linewidth}}{
        \cellcolor[rgb]{0.9,0.9,0.9}{
            \makecell[{{p{\linewidth}}}]{
                \texttt{\tiny{[GM$|$GM]}}
                \texttt{BEWEEG: P, (3, 6)} \\
            }
        }
    }
    & & \\ \\

    \theutterance \stepcounter{utterance}  
    & & & \multicolumn{2}{p{0.3\linewidth}}{
        \cellcolor[rgb]{0.9,0.9,0.9}{
            \makecell[{{p{\linewidth}}}]{
                \texttt{\tiny{[GM$|$GM]}}
                \texttt{Straf: (3,6) is niet leeg, maar bevat '─'.} \\
\\ 
\texttt{Jullie hebben gezamenlijk 10 van 8 straffen verzameld.} \\
\texttt{Probeer het opnieuw!} \\
            }
        }
    }
    & & \\ \\

    \theutterance \stepcounter{utterance}  
    & & \multicolumn{4}{p{0.6\linewidth}}{
        \cellcolor[rgb]{0.9,0.9,0.9}{
            \makecell[{{p{\linewidth}}}]{
                \texttt{\tiny{[GM$\rangle$P2]}}
                \texttt{Straf: (3,6) is niet leeg, maar bevat '─'.} \\
\\ 
\texttt{Jullie hebben gezamenlijk 10 van 8 straffen verzameld.} \\
\texttt{Probeer het opnieuw!} \\
            }
        }
    }
    & \\ \\

    \theutterance \stepcounter{utterance}  
    & & & \multicolumn{4}{p{0.6\linewidth}}{
        \cellcolor[rgb]{0.9,0.9,0.9}{
            \makecell[{{p{\linewidth}}}]{
                \texttt{\tiny{[GM$\langle$P2]}}
                \texttt{BEWEEG: P, (3, 5)} \\
            }
        }
    }
     \\ \\

    \theutterance \stepcounter{utterance}  
    & & & \multicolumn{2}{p{0.3\linewidth}}{
        \cellcolor[rgb]{0.9,0.9,0.9}{
            \makecell[{{p{\linewidth}}}]{
                \texttt{\tiny{[GM$|$GM]}}
                \texttt{BEWEEG: P, (3, 5)} \\
            }
        }
    }
    & & \\ \\

    \theutterance \stepcounter{utterance}  
    & & & \multicolumn{2}{p{0.3\linewidth}}{
        \cellcolor[rgb]{0.9,0.9,0.9}{
            \makecell[{{p{\linewidth}}}]{
                \texttt{\tiny{[GM$|$GM]}}
                \texttt{P moved to (3, 5)} \\
            }
        }
    }
    & & \\ \\

    \theutterance \stepcounter{utterance}  
    & & \multicolumn{4}{p{0.6\linewidth}}{
        \cellcolor[rgb]{0.9,0.9,0.9}{
            \makecell[{{p{\linewidth}}}]{
                \texttt{\tiny{[P1$\langle$GM]}}
                \texttt{Je bericht is doorgestuurd naar de andere speler.} \\
\\ 
\\ 
\texttt{Je speelt momenteel ronde 9 van maximaal 12.} \\
\texttt{Jullie hebben gezamenlijk 10 van 8 straffen verzameld.} \\
\texttt{De andere speler heeft een voorwerp op zijn rooster verplaatst.} \\
\\ 
\\ 
\texttt{Wat is je volgende opdracht?} \\
            }
        }
    }
    & \\ \\

    \theutterance \stepcounter{utterance}  
    & \multicolumn{4}{p{0.6\linewidth}}{
        \cellcolor[rgb]{0.9,0.9,0.9}{
            \makecell[{{p{\linewidth}}}]{
                \texttt{\tiny{[P1$\rangle$GM]}}
                \texttt{ZEG: Oei, we hebben de 8 straffen overschreden. Dat betekent dat we het spel hebben verloren. Ik hoop dat de andere speler nog een bericht heeft gestuurd over zijn laatste zet.} \\
            }
        }
    }
    & & \\ \\

    \theutterance \stepcounter{utterance}  
    & & & \multicolumn{2}{p{0.3\linewidth}}{
        \cellcolor[rgb]{0.9,0.9,0.9}{
            \makecell[{{p{\linewidth}}}]{
                \texttt{\tiny{[GM$|$GM]}}
                \texttt{ZEG: Oei, we hebben de 8 straffen overschreden. Dat betekent dat we het spel hebben verloren. Ik hoop dat de andere speler nog een bericht heeft gestuurd over zijn laatste zet.} \\
            }
        }
    }
    & & \\ \\

    \theutterance \stepcounter{utterance}  
    & & \multicolumn{4}{p{0.6\linewidth}}{
        \cellcolor[rgb]{0.9,0.9,0.9}{
            \makecell[{{p{\linewidth}}}]{
                \texttt{\tiny{[GM$\rangle$P2]}}
                \texttt{Verplaatst 'P' naar (3,5) met succes. Uw bijgewerkte raster ziet er als volgt uit:} \\
\texttt{```} \\
\texttt{ 1234567} \\
\texttt{╔═╤══╤══╗} \\
\texttt{║◌│◌◌│◌◌║ 1} \\
\texttt{║◌└┬─┴──╢ 2} \\
\texttt{║◌◌│◌◌◌◌║ 3} \\
\texttt{╟─┬┘L◌┌─╢ 4} \\
\texttt{║◌│PC◌│◌║ 5} \\
\texttt{╟─┼─┬─┤◌║ 6} \\
\texttt{║◌│◌│◌│◌║ 7} \\
\texttt{╚═╧═╧═╧═╝} \\
\texttt{```} \\
\texttt{Je speelt momenteel ronde 10 van maximaal 12.} \\
\texttt{Jullie hebben gezamenlijk 10 van 8 straffen verzameld.} \\
\texttt{De andere speler stuurde dit bericht:} \\
\\ 
\\ 
\texttt{"Oei, we hebben de 8 straffen overschreden. Dat betekent dat we het spel hebben verloren. Ik hoop dat de andere speler nog een bericht heeft gestuurd over zijn laatste zet."} \\
\\ 
\\ 
\texttt{Wat is je volgende opdracht?} \\
            }
        }
    }
    & \\ \\

    \theutterance \stepcounter{utterance}  
    & & & \multicolumn{4}{p{0.6\linewidth}}{
        \cellcolor[rgb]{0.9,0.9,0.9}{
            \makecell[{{p{\linewidth}}}]{
                \texttt{\tiny{[GM$\langle$P2]}}
                \texttt{ZEG: Het spijt me, ik heb de fouten niet kunnen voorkomen. We hebben het spel verloren.} \\
            }
        }
    }
     \\ \\

    \theutterance \stepcounter{utterance}  
    & & & \multicolumn{2}{p{0.3\linewidth}}{
        \cellcolor[rgb]{0.9,0.9,0.9}{
            \makecell[{{p{\linewidth}}}]{
                \texttt{\tiny{[GM$|$GM]}}
                \texttt{ZEG: Het spijt me, ik heb de fouten niet kunnen voorkomen. We hebben het spel verloren.} \\
            }
        }
    }
    & & \\ \\

    \theutterance \stepcounter{utterance}  
    & & \multicolumn{4}{p{0.6\linewidth}}{
        \cellcolor[rgb]{0.9,0.9,0.9}{
            \makecell[{{p{\linewidth}}}]{
                \texttt{\tiny{[P1$\langle$GM]}}
                \texttt{Je bericht is doorgestuurd naar de andere speler.} \\
\\ 
\\ 
\texttt{Je speelt momenteel ronde 10 van maximaal 12.} \\
\texttt{Jullie hebben gezamenlijk 10 van 8 straffen verzameld.} \\
\texttt{De andere speler stuurde dit bericht:} \\
\\ 
\\ 
\texttt{"Het spijt me, ik heb de fouten niet kunnen voorkomen. We hebben het spel verloren."} \\
\\ 
\\ 
\texttt{Wat is je volgende opdracht?} \\
            }
        }
    }
    & \\ \\

    \theutterance \stepcounter{utterance}  
    & \multicolumn{4}{p{0.6\linewidth}}{
        \cellcolor[rgb]{0.9,0.9,0.9}{
            \makecell[{{p{\linewidth}}}]{
                \texttt{\tiny{[P1$\rangle$GM]}}
                \texttt{ZEG: Het spijt me ook, het was een leerzame ervaring. We hadden net een goede strategie, maar de uitvoering liep helaas mis door de strafpunten. Bedankt voor het samen spelen!} \\
            }
        }
    }
    & & \\ \\

    \theutterance \stepcounter{utterance}  
    & & & \multicolumn{2}{p{0.3\linewidth}}{
        \cellcolor[rgb]{0.9,0.9,0.9}{
            \makecell[{{p{\linewidth}}}]{
                \texttt{\tiny{[GM$|$GM]}}
                \texttt{ZEG: Het spijt me ook, het was een leerzame ervaring. We hadden net een goede strategie, maar de uitvoering liep helaas mis door de strafpunten. Bedankt voor het samen spelen!} \\
            }
        }
    }
    & & \\ \\

    \theutterance \stepcounter{utterance}  
    & & \multicolumn{4}{p{0.6\linewidth}}{
        \cellcolor[rgb]{0.9,0.9,0.9}{
            \makecell[{{p{\linewidth}}}]{
                \texttt{\tiny{[GM$\rangle$P2]}}
                \texttt{Je bericht is doorgestuurd naar de andere speler.} \\
\\ 
\\ 
\texttt{Je speelt momenteel ronde 11 van maximaal 12.} \\
\texttt{Jullie hebben gezamenlijk 10 van 8 straffen verzameld.} \\
\texttt{De andere speler stuurde dit bericht:} \\
\\ 
\\ 
\texttt{"Het spijt me ook, het was een leerzame ervaring. We hadden net een goede strategie, maar de uitvoering liep helaas mis door de strafpunten. Bedankt voor het samen spelen!"} \\
\\ 
\\ 
\texttt{Wat is je volgende opdracht?} \\
            }
        }
    }
    & \\ \\

    \theutterance \stepcounter{utterance}  
    & & & \multicolumn{4}{p{0.6\linewidth}}{
        \cellcolor[rgb]{0.9,0.9,0.9}{
            \makecell[{{p{\linewidth}}}]{
                \texttt{\tiny{[GM$\langle$P2]}}
                \texttt{ZEG: Bedankt voor het samen spelen! Het was een goede poging en ik heb er ook van geleerd. Tot de volgende keer!} \\
            }
        }
    }
     \\ \\

    \theutterance \stepcounter{utterance}  
    & & & \multicolumn{2}{p{0.3\linewidth}}{
        \cellcolor[rgb]{0.9,0.9,0.9}{
            \makecell[{{p{\linewidth}}}]{
                \texttt{\tiny{[GM$|$GM]}}
                \texttt{ZEG: Bedankt voor het samen spelen! Het was een goede poging en ik heb er ook van geleerd. Tot de volgende keer!} \\
            }
        }
    }
    & & \\ \\

    \theutterance \stepcounter{utterance}  
    & & \multicolumn{4}{p{0.6\linewidth}}{
        \cellcolor[rgb]{0.9,0.9,0.9}{
            \makecell[{{p{\linewidth}}}]{
                \texttt{\tiny{[P1$\langle$GM]}}
                \texttt{Je bericht is doorgestuurd naar de andere speler.} \\
\\ 
\\ 
\texttt{Je speelt momenteel ronde 11 van maximaal 12.} \\
\texttt{Jullie hebben gezamenlijk 10 van 8 straffen verzameld.} \\
\texttt{De andere speler stuurde dit bericht:} \\
\\ 
\\ 
\texttt{"Bedankt voor het samen spelen! Het was een goede poging en ik heb er ook van geleerd. Tot de volgende keer!"} \\
\\ 
\\ 
\texttt{Wat is je volgende opdracht?} \\
            }
        }
    }
    & \\ \\

    \theutterance \stepcounter{utterance}  
    & \multicolumn{4}{p{0.6\linewidth}}{
        \cellcolor[rgb]{0.9,0.9,0.9}{
            \makecell[{{p{\linewidth}}}]{
                \texttt{\tiny{[P1$\rangle$GM]}}
                \texttt{ZEG: Tot de volgende keer!} \\
            }
        }
    }
    & & \\ \\

    \theutterance \stepcounter{utterance}  
    & & & \multicolumn{2}{p{0.3\linewidth}}{
        \cellcolor[rgb]{0.9,0.9,0.9}{
            \makecell[{{p{\linewidth}}}]{
                \texttt{\tiny{[GM$|$GM]}}
                \texttt{ZEG: Tot de volgende keer!} \\
            }
        }
    }
    & & \\ \\

    \theutterance \stepcounter{utterance}  
    & & \multicolumn{4}{p{0.6\linewidth}}{
        \cellcolor[rgb]{0.9,0.9,0.9}{
            \makecell[{{p{\linewidth}}}]{
                \texttt{\tiny{[GM$\rangle$P2]}}
                \texttt{Je bericht is doorgestuurd naar de andere speler.} \\
\\ 
\\ 
\texttt{Je speelt momenteel ronde 12 van maximaal 12.} \\
\texttt{Jullie hebben gezamenlijk 10 van 8 straffen verzameld.} \\
\texttt{De andere speler stuurde dit bericht:} \\
\\ 
\\ 
\texttt{"Tot de volgende keer!"} \\
\\ 
\\ 
\texttt{Wat is je volgende opdracht?} \\
            }
        }
    }
    & \\ \\

    \theutterance \stepcounter{utterance}  
    & & & \multicolumn{4}{p{0.6\linewidth}}{
        \cellcolor[rgb]{0.9,0.9,0.9}{
            \makecell[{{p{\linewidth}}}]{
                \texttt{\tiny{[GM$\langle$P2]}}
                \texttt{ZEG: Tot ziens!} \\
            }
        }
    }
     \\ \\

    \theutterance \stepcounter{utterance}  
    & & & \multicolumn{2}{p{0.3\linewidth}}{
        \cellcolor[rgb]{0.9,0.9,0.9}{
            \makecell[{{p{\linewidth}}}]{
                \texttt{\tiny{[GM$|$GM]}}
                \texttt{ZEG: Tot ziens!} \\
            }
        }
    }
    & & \\ \\

    \theutterance \stepcounter{utterance}  
    & & & \multicolumn{2}{p{0.3\linewidth}}{
        \cellcolor[rgb]{0.9,0.9,0.9}{
            \makecell[{{p{\linewidth}}}]{
                \texttt{\tiny{[GM$|$GM]}}
                \texttt{Maximum number of rounds reached} \\
            }
        }
    }
    & & \\ \\

    \theutterance \stepcounter{utterance}  
    & & & \multicolumn{2}{p{0.3\linewidth}}{
        \cellcolor[rgb]{0.9,0.9,0.9}{
            \makecell[{{p{\linewidth}}}]{
                \texttt{\tiny{[GM$|$GM]}}
                \texttt{Initial states:} \\
\texttt{```} \\
\texttt{Player 1:} \\
\texttt{ 1234567} \\
\texttt{╔═╤══╤══╗} \\
\texttt{║◌│◌◌│◌◌║ 1} \\
\texttt{║◌└┬─┴──╢ 2} \\
\texttt{║◌◌│L◌◌◌║ 3} \\
\texttt{╟─┬┘◌◌┌─╢ 4} \\
\texttt{║◌│◌◌◌│C║ 5} \\
\texttt{╟─┼─┬─┤P║ 6} \\
\texttt{║◌│◌│◌│◌║ 7} \\
\texttt{╚═╧═╧═╧═╝} \\
\\ 
\\ 
\texttt{Player 2:} \\
\texttt{ 1234567} \\
\texttt{╔═╤══╤══╗} \\
\texttt{║◌│◌◌│◌◌║ 1} \\
\texttt{║L└┬─┴──╢ 2} \\
\texttt{║◌C│◌◌◌◌║ 3} \\
\texttt{╟─┬┘◌◌┌─╢ 4} \\
\texttt{║◌│◌◌◌│◌║ 5} \\
\texttt{╟─┼─┬─┤◌║ 6} \\
\texttt{║◌│P│◌│◌║ 7} \\
\texttt{╚═╧═╧═╧═╝} \\
\\ 
\texttt{```} \\
            }
        }
    }
    & & \\ \\

    \theutterance \stepcounter{utterance}  
    & & & \multicolumn{2}{p{0.3\linewidth}}{
        \cellcolor[rgb]{0.9,0.9,0.9}{
            \makecell[{{p{\linewidth}}}]{
                \texttt{\tiny{[GM$|$GM]}}
                \texttt{End states:} \\
\texttt{```} \\
\texttt{Player 1:} \\
\texttt{ 1234567} \\
\texttt{╔═╤══╤══╗} \\
\texttt{║◌│◌◌│◌◌║ 1} \\
\texttt{║◌└┬─┴──╢ 2} \\
\texttt{║◌◌│◌◌◌◌║ 3} \\
\texttt{╟─┬┘L◌┌─╢ 4} \\
\texttt{║◌│◌C◌│◌║ 5} \\
\texttt{╟─┼─┬─┤◌║ 6} \\
\texttt{║◌│◌│P│◌║ 7} \\
\texttt{╚═╧═╧═╧═╝} \\
\\ 
\\ 
\texttt{Player 2:} \\
\texttt{ 1234567} \\
\texttt{╔═╤══╤══╗} \\
\texttt{║◌│◌◌│◌◌║ 1} \\
\texttt{║◌└┬─┴──╢ 2} \\
\texttt{║◌◌│◌◌◌◌║ 3} \\
\texttt{╟─┬┘L◌┌─╢ 4} \\
\texttt{║◌│PC◌│◌║ 5} \\
\texttt{╟─┼─┬─┤◌║ 6} \\
\texttt{║◌│◌│◌│◌║ 7} \\
\texttt{╚═╧═╧═╧═╝} \\
\\ 
\texttt{```} \\
            }
        }
    }
    & & \\ \\

    \theutterance \stepcounter{utterance}  
    & & & \multicolumn{2}{p{0.3\linewidth}}{
        \cellcolor[rgb]{0.9,0.9,0.9}{
            \makecell[{{p{\linewidth}}}]{
                \texttt{\tiny{[GM$|$GM]}}
                \texttt{* success: True} \\
\texttt{* lose: False} \\
\texttt{* aborted: False} \\
\texttt{{-}{-}{-}{-}{-}{-}{-}} \\
\texttt{* Shifts: 2.00} \\
\texttt{* Max Shifts: 4.00} \\
\texttt{* Min Shifts: 2.00} \\
\texttt{* End Distance Sum: 2.83} \\
\texttt{* Init Distance Sum: 12.67} \\
\texttt{* Expected Distance Sum: 12.57} \\
\texttt{* Penalties: 10.00} \\
\texttt{* Max Penalties: 8.00} \\
\texttt{* Rounds: 12.00} \\
\texttt{* Max Rounds: 12.00} \\
\texttt{* Object Count: 3.00} \\
            }
        }
    }
    & & \\ \\

\end{supertabular}
}

\end{document}
